\emph{Översikt}
Den här modulen är mötet med Python som språk.  En algoritm kan förfinas
ända ner till körbar kod, som vi gjorde i föreläsningen \emph{Algoritmiskt
tänkande}; nu frågar vi vad datorn gör med den koden.  Svaret är enkelt
att säga och svårt att vänja sig vid: den läser programmet \emph{uppifrån
och ned} och exekverar \emph{en sats i taget}.  \emph{Variabler} är namn
som håller ett värde; en \emph{tilldelning} beräknar högerledet ur de
värden som gäller just då och ersätter det gamla värdet.  Värdenas
\emph{typer} --- heltal, flyttal och strängar --- avgör vilka operationer
som fungerar.  Vi skriver ut värden med \mintinline{python}{print} och
f-strängar och avslutar med det som gör kod läsbar för människor:
beskrivande namn, konstanter, kommentarer och stilreglerna i PEP~8.

\emph{Lärandemål}
Efter denna modul ska studenten kunna:
% Lo03
\begin{restatable}{lo}{VarLOExecution}\label{VarLOExecution}%
  Redogöra för vad som händer när ett pythonprogram körs, och följa
  exekveringen rad för rad --- med papper och penna eller med ett
  spårningsverktyg som Python Tutor.
\end{restatable}
% Lo03
\begin{restatable}{lo}{VarLOAssignment}\label{VarLOAssignment}%
  Lagra värden i variabler och förklara vad en tilldelning gör.
\end{restatable}
% Lo03
\begin{restatable}{lo}{VarLOTypes}\label{VarLOTypes}%
  Skilja på datatyperna \mintinline{python}{int}, \mintinline{python}{float}
  och \mintinline{python}{str}, avgöra vilka operationer som passar för
  vilken typ, och konvertera mellan dem.
\end{restatable}
% Lo05
\begin{restatable}{lo}{VarLOOutput}\label{VarLOOutput}%
  Skriva ut text och värden med \mintinline{python}{print} och f-strängar
  så att utskriften blir läsbar för den som kör programmet.
\end{restatable}
% Lo07
\begin{restatable}{lo}{VarLONames}\label{VarLONames}%
  Välja beskrivande variabelnamn och markera konstanter enligt PEP~8, och
  förklara vem namnet är till för.
\end{restatable}
% Lo07
\begin{restatable}{lo}{VarLOComments}\label{VarLOComments}%
  Kommentera kod så att kommentaren förklarar \emph{varför}, inte
  \emph{vad}, och skriva kod som följer PEP~8.
\end{restatable}

\emph{Förkunskaper}
Studenten bör ha tagit del av föreläsningarna \emph{Hello World!} och
\emph{Algoritmiskt tänkande}: hon ska kunna köra ett pythonprogram från
terminalen, redigera det i en textredigerare, läsa notationen med
namngivna steg och utveckla en algoritm genom stegvis förfining.

\emph{Läsning}
Kompletterande material finns i kursens videogenomgång \emph{Variabler och
utskrifter}.  Stilreglerna som modulen hänvisar till står i
PEP~8\autocite{pep8}.
